\documentclass[10pt,a4paper]{article}

\usepackage[utf8]{inputenc}
\usepackage[T1]{fontenc}
\usepackage[french]{babel}
\usepackage{amsmath,amssymb}
\usepackage{geometry}
\usepackage{enumitem}
\usepackage{multicol}
\usepackage{xcolor}
\usepackage{mdframed}
\usepackage{lmodern}
\usepackage{microtype}
\usepackage{fancyhdr}
\usepackage{titlesec}
\usepackage{parskip}

\geometry{top=1.8cm, bottom=2cm, left=1.8cm, right=1.8cm}

\pagestyle{fancy}
\fancyhf{}
\fancyhead[L]{\small\textit{Mathématiques, fiche de révision}}
\fancyhead[R]{\small\textit{Fonctions de plusieurs variables}}
\fancyfoot[C]{\small\thepage}
\fancyfoot[R]{\tiny\textcopyright\ Rohan Fossé}
\renewcommand{\headrulewidth}{0.3pt}

\definecolor{bleu}{RGB}{30,60,120}
\definecolor{vert}{RGB}{0,110,60}
\definecolor{rouge}{RGB}{160,20,20}
\definecolor{gris}{gray}{0.93}
\definecolor{jauneclair}{RGB}{255,250,220}
\definecolor{vertclair}{RGB}{220,245,220}

\newmdenv[backgroundcolor=gris, linewidth=0pt,
  innerleftmargin=8pt, innerrightmargin=8pt,
  innertopmargin=5pt, innerbottommargin=5pt,
  skipabove=4pt, skipbelow=4pt]{grisbox}

\newmdenv[backgroundcolor=jauneclair, linecolor=rouge, linewidth=1pt,
  innerleftmargin=8pt, innerrightmargin=8pt,
  innertopmargin=5pt, innerbottommargin=5pt,
  skipabove=4pt, skipbelow=4pt]{alertbox}

\titleformat{\section}[block]{\normalsize\bfseries\color{bleu}}{}{0pt}{}[\vspace{-3pt}\rule{\linewidth}{0.5pt}]
\titlespacing{\section}{0pt}{10pt}{4pt}

\titleformat{\subsection}[runin]{\small\bfseries\color{vert}}{}{0pt}{}[\ :]
\titlespacing{\subsection}{0pt}{4pt}{0pt}

\begin{document}

\begin{center}
  {\Large\bfseries\color{bleu} Fiche de révision : Fonctions de plusieurs variables}\\[3pt]
  {\small Continuité · Différentielle · Dérivées partielles · Plan tangent · Extrema}
\end{center}
\noindent\rule{\linewidth}{0.8pt}

\begin{multicols}{2}

\section{Continuité}

\begin{grisbox}
\textbf{Définition.} $f:\mathbb{R}^2\to\mathbb{R}$ est \textbf{continue en $a=(a_1,a_2)$} si :
$$\lim_{(x,y)\to a} f(x,y) = f(a)$$
\end{grisbox}

\begin{alertbox}
\textbf{Attention !} L'existence des dérivées partielles ne suffit \textit{pas} à garantir la continuité (ni la différentiabilité).
\end{alertbox}

\textbf{Tester la non-existence d'une limite} : tester deux chemins différents aboutissant à des valeurs distinctes.

\section{Dérivées partielles}

\begin{grisbox}
\textbf{Définition.}
$$\frac{\partial f}{\partial x}(a,b) = \lim_{h\to 0}\frac{f(a+h,b)-f(a,b)}{h}$$
On dérive par rapport à $x$ en traitant $y$ comme une \textbf{constante}.
\end{grisbox}

\subsection{Règle de la chaîne}
Si $f(x,y)=g(h(x,y))$, alors :
$$\frac{\partial f}{\partial x} = g'(h)\cdot\frac{\partial h}{\partial x}$$

\subsection{Théorème de Schwartz}
Si $f\in\mathcal{C}^2$, alors les dérivées croisées sont égales :
$$\boxed{\frac{\partial^2 f}{\partial x\partial y} = \frac{\partial^2 f}{\partial y\partial x}}$$

\section{Différentielle totale}

\begin{grisbox}
\textbf{Définition.} Si $f$ est différentiable en $(a,b)$ :
$$df = \frac{\partial f}{\partial x}\,dx + \frac{\partial f}{\partial y}\,dy$$
La différentiabilité \textbf{implique la continuité}.
\end{grisbox}

\section{Plan tangent}

\begin{grisbox}
\textbf{Surface $z=f(x,y)$} au point $(a,b,f(a,b))$ :
$$\boxed{z = f(a,b) + \frac{\partial f}{\partial x}(a,b)(x-a) + \frac{\partial f}{\partial y}(a,b)(y-b)}$$

\textbf{Surface implicite $F(x,y,z)=0$} au point $M_0$ :
$$\nabla F(M_0)\cdot(x-x_0,y-y_0,z-z_0)=0$$
Le vecteur normal est $\nabla F(M_0)$.
\end{grisbox}

\section{Extrema locaux}

\begin{grisbox}
\textbf{Étape 1 -- Points critiques :} résoudre $\nabla f = \mathbf{0}$
$$\frac{\partial f}{\partial x}=0 \quad \text{et} \quad \frac{\partial f}{\partial y}=0$$

\textbf{Étape 2 -- Discriminant hessien :}
$$H = f_{xx}\,f_{yy} - (f_{xy})^2$$

\textbf{Étape 3 -- Conclusion :}
\begin{itemize}[leftmargin=1em,itemsep=1pt]
  \item $H > 0$ et $f_{xx}>0$ : \textbf{minimum local}
  \item $H > 0$ et $f_{xx}<0$ : \textbf{maximum local}
  \item $H < 0$ : \textbf{point selle} (col)
  \item $H = 0$ : pas de conclusion
\end{itemize}
\end{grisbox}

\section{À retenir absolument}

\begin{grisbox}
\begin{itemize}[leftmargin=1em,itemsep=2pt]
  \item \textbf{Continuité} : $\lim = f(a)$
  \item \textbf{$\partial f/\partial x$} : dériver en $x$, $y$ constant
  \item \textbf{Schwartz} : $f_{xy}=f_{yx}$ si $f\in\mathcal{C}^2$
  \item \textbf{Différentielle} implique continuité
  \item \textbf{Plan tangent} : formule DL d'ordre 1
  \item \textbf{Point critique} : $\nabla f=\mathbf{0}$
  \item \textbf{Nature} : signe de $H$ et $f_{xx}$
\end{itemize}
\end{grisbox}

\end{multicols}

\end{document}
